% =====================================================================
%  19.4  The Iterative (Bottom-Up) Segment Tree       sec:c13-iter
%  Exports: prop:c13-iter  lst:c13-iter  lst:c13-iter-range
%           lst:c13-514-soln  fig:c13-iter2n  prob:c13-r2d2
%  Solved: CF 514D R2D2 and Droid Army (2000), live-verified.
%    Statement: n droids in a row, each has m (<=5) detail counts; m weapons,
%    weapon j must fire (max over the window of column j) shots to clear type j
%    across a consecutive run; at most k shots total. Maximise the run length;
%    print the m per-column maxima of ONE such optimal run (any if ties; if the
%    best run is empty, print m zeros).  Output is ONE line of m integers.
%    Constraints (verified via DOM <sup>): 1<=n<=1e5, 1<=m<=5, 0<=k<=1e9,
%    0<=a_i<=1e8.  (innerText had collapsed 1e5/1e9/1e8 to 105/109/108.)
%    Samples (verbatim from CF):
%      1) "5 2 4 / 4 0 / 1 2 / 2 1 / 0 2 / 1 3"  -> "2 2"
%         (destroy droids 2,3,4; per-column max = (max 1,2,0, max 2,1,2)=(2,2),
%          sum 4<=k; length 3 is maximal -- droids1-3 cost 6, droids2-5 cost 5)
%      2) "3 2 4 / 1 2 / 1 3 / 2 2"              -> "1 3"
%         (destroy droids 1,2; per-column max=(1,3), sum 4<=k; length-3 cost 5)
%  Running numbers agree with 00_opener (fig:c13-tree): the 2n build over the
%    shared 8-array (3,2,4,5,1,1,5,3) re-uses 19.1's min-tree t[1..15] exactly,
%    since n=8 is a power of two -- the two layouts differ only in the walk.
%  Discharges Ch.18's forward pointer: \secref{c14-range} drew the BIT-vs-segtree
%    *expressiveness* line; this section settles the remaining axes (memory,
%    constant factor, code, and the half-open walk), NOT re-listing expressiveness.
% =====================================================================
\section{The Iterative (Bottom-Up) Segment Tree}\label{sec:c13-iter}

\begin{previouslybox}[Ch.~18 lived in one array with a \text{lowbit}; so can we]
Chapter~18's entire power came from throwing away the tree \emph{shape} and
keeping only an array indexed by the binary carries of the position.  A segment
tree is usually presented the opposite way -- recursion, node intervals, three
functions that call themselves.  But when the length is fixed, the tree is a
\emph{perfect} binary tree, and a perfect binary tree fits in a flat array with
no interval arguments at all: leaves in the back half, internal nodes in the
front half, parent \(= i/2\), sibling \(= i \oplus 1\).  Everything in
\(\secref{c13-core}\) survives; only the recursion is deleted.  This section
builds that \(2n\) form, gives its half-open range walk, and then does the
honest accounting Ch.~18 asked for: what the iterative tree buys, what it costs,
and exactly when the recursive skeleton of \(\secref{c13-core}\) is the better
tool.
\end{previouslybox}

\begin{topicintro}
The recursive tree of \(\secref{c13-core}\) carries a node's interval \([tl,tr]\)
as \emph{arguments}, because the \(4n\) heap layout lets the tree be ragged when
\(n\) is not a power of two.  Pad the array to the next power of two (or, for a
commutative fold, skip padding entirely) and the tree becomes complete: all
leaves sit at depth \(\lceil\log_2 n\rceil\), and a complete binary tree over \(n\)
leaves occupies exactly \(2n\) array slots, indices \(1..2n-1\), with the leaves
at \(t[n..2n{-}1]\) and the internal nodes at \(t[1..n{-}1]\).  No recursion, no
interval parameters, no \(4n\) slack -- just two loops and one walk.

\textbf{Prerequisites.} \(\secref{c13-core}\) (heap layout, the fold \(\mathcal{M}\),
the identity \(e\)).  \(\secref{c13-combine}\) (a node may be a record; order can
matter).  The opener's \(2n\) fold table over the shared array \(A\)
(\(\figref{c13-tree}\)).  Chapter~18's \(\secref{c14-core}\) for the array-only
mindset this mirrors.

\textbf{After this section you can}
\begin{itemize}[nosep]
  \item write the \(2n\) build, point update, and half-open query loops from
        memory, and say which fold each requires to be commutative; \hfill\lv{3}--\lv{4}
  \item explain, with the two-accumulator walk, why a non-commutative combine is
        where the recursive \(\coderef{c13-query}\) stays simpler; \hfill\lv{4}--\lv{5}
  \item fill in the honest BIT / recursive-segtree / iterative-segtree cost table
        that Ch.~18 left open; \hfill\lv{3}
  \item solve CF~514D end to end -- \(m\) iterative max-trees plus a binary search
        on the answer. \hfill\lv{5}
\end{itemize}
\end{topicintro}

\begin{recognitionbox}[Pattern card -- ``fixed length, commutative fold, want terse code'']
\reckey{the array size never changes across the whole run; the fold is
\(\min/\max/\text{sum}/\gcd/\text{xor}/\)OR/AND (all commutative); you want the
smallest constant factor and the least code; a range query over a
\emph{half-open} \([l,r)\) window; or a problem that pairs a segment tree with an
outer binary search so the \(\log n\) is spent inside a loop and must be cheap.}
\recuse{
  \item static length + commutative fold -- this \(2n\) form, it is the fastest to
        type and runs;
  \item the fold is non-commutative but associative (max-subarray, brackets) --
        still fine here \emph{if} you keep both accumulators, or fall back to the
        recursive \(\coderef{c13-query}\);
  \item the tree is an inner subroutine called \(\bigO{\log}\) times -- the low
        constant matters, use this form.
}
\recnotuse{
  \item you must \emph{descend} for a threshold (\(\secref{c13-search}\)) -- the
        bottom-up layout has no clean interval at each node, so the recursive
        descend (or a padded power-of-two variant) reads better;
  \item the array grows or nodes are inserted -- a balanced/BST structure
        (Ch.~6) or a sparse dynamic tree, not a fixed \(2n\) array;
  \item range \emph{updates} -- lazy tags need the top-down interval, Ch.~20
        (\(\chref{c15}\)); this section keeps every write a single point.
}
One-line test: \emph{does the length stay put and is your combine commutative?
Then the \(2n\) array-resident tree is the leanest thing that folds a range.}
\end{recognitionbox}

% ---------------------------------------------------------------
\subsection{The \(2n\) Layout and Its Two Loops}
% ---------------------------------------------------------------

Put the leaves in the second half of the array and index a node's children the
same way the recursive heap layout does -- \(2i\) and \(2i{+}1\) -- and a
complete tree needs no arguments: a node's parent is just \(i/2\), so an update
climbs by halving.

\begin{figure}[H]
\centering
\begin{tikzpicture}[diagram, scale=0.92, transform shape]
  % leaves t[8..15]
  \foreach \i/\vv in {0/3,1/2,2/4,3/5,4/1,5/1,6/5,7/3} {
    \draw[primaryblue, line width=0.5pt, fill=primaryblue!6] (\i*1.0,0) rectangle (\i*1.0+0.9,0.7);
    \node[font=\scriptsize\ttfamily] at (\i*1.0+0.45,0.35) {\vv};
    \node[font=\tiny\ttfamily, text=darkgray] at (\i*1.0+0.45,-0.26) {\(t_{\the\numexpr8+\i\relax}\)};}
  % internal t[4..7]
  \foreach \i/\x/\val in {4/0.45/2,5/2.45/4,6/4.45/1,7/6.45/3} {
    \node[boxnode, successgreen, minimum size=6mm, font=\scriptsize] (n\i) at (\x,1.5) {\(t_{\i}{=}\val\)};}
  % internal t[2..3]
  \node[boxnode, successgreen, minimum size=6mm, font=\scriptsize] (n2) at (1.45,2.9) {\(t_{2}{=}2\)};
  \node[boxnode, successgreen, minimum size=6mm, font=\scriptsize] (n3) at (5.45,2.9) {\(t_{3}{=}1\)};
  % root
  \node[boxnode, successgreen, minimum size=7mm, line width=1pt, fill=successgreen!16, font=\scriptsize] (n1) at (3.45,4.2) {\(t_{1}{=}1\)};
  % edges
  \draw[flow] (n1)--(n2); \draw[flow] (n1)--(n3);
  \draw[flow] (n2)--(n4); \draw[flow] (n2)--(n5);
  \draw[flow] (n3)--(n6); \draw[flow] (n3)--(n7);
  \draw[flow] (n4)--(0.45,0.7); \draw[flow] (n4)--(1.45,0.7);
  \draw[flow] (n5)--(2.45,0.7); \draw[flow] (n5)--(3.45,0.7);
  \draw[flow] (n6)--(4.45,0.7); \draw[flow] (n6)--(5.45,0.7);
  \draw[flow] (n7)--(6.45,0.7); \draw[flow] (n7)--(7.45,0.7);
  \node[annot, font=\scriptsize, text=primaryblue, anchor=east] at (-0.1,0.35) {leaves \(t_{8..15}{=}A\)};
  \node[annot, font=\scriptsize, text=green!45!black, anchor=east] at (-0.1,4.2) {internal \(t_{1..7}\)};
\end{tikzpicture}
\caption{The \(2n\) array-resident tree over the same \(n{=}8\) min-instance
\((3,2,4,5,1,1,5,3)\) as \(\figref{c13-minbuild}\).  Leaves occupy the back half
\(t[8..15]\); the front half \(t[1..7]\) is filled by \(\coderef{c13-iter}\)'s
second loop, \(t_i=\min(t_{2i},t_{2i+1})\), from \(i{=}7\) down to \(1\).  The
values are \emph{identical} to the recursive \(4n\) tree because \(n{=}8\) is a
power of two -- only the walk changes, and no slots are wasted.}
\label{fig:c13-iter2n}
\end{figure}

\begin{lstlisting}[caption={the 2n bottom-up tree: build and point update, no recursion}, label={lst:c13-iter}]
int n;                        // number of leaves, fixed for the whole run
vector<long long> t;          // 2*n slots, indices 1 .. 2n-1 used
long long mergeOp(long long a, long long b){ return min(a, b); }  // the fold M
const long long ID = (1LL<<60);   // identity e: min->+inf, sum/gcd->0, max->-inf

void build(const vector<long long>& a){      // a indexed 0..n-1
    t.assign(2*n, ID);
    for (int i = 0; i < n; i++) t[n+i] = a[i];        // leaves in place
    for (int i = n-1; i > 0; i--) t[i] = mergeOp(t[2*i], t[2*i+1]);  // bottom-up
}
// O(n): one merge per internal node, exactly like the recursive build
void update(int p, long long val){           // set leaf p (0-based) to val
    for (t[p += n] = val; p > 1; p >>= 1)    // climb to the root: parent = p/2
        t[p >> 1] = mergeOp(t[p], t[p ^ 1]); // p^1 is p's sibling
}
// O(log n): touches only the ancestors of leaf p
\end{lstlisting}

\noindent The sibling trick \code{p\textasciicircum1} is the whole point of the
array layout: because leaves and internal nodes share one heap, a node at index
\(p\) and its sibling are the pair \((p, p\oplus 1)\) -- even index gets its odd
neighbour and vice versa -- so the parent \(p/2\) can be re-folded from the just
updated child and its untouched sibling without ever storing intervals.

% ---------------------------------------------------------------
\subsection{The Half-Open Range Query}
% ---------------------------------------------------------------

The recursive query \(\coderef{c13-query}\) passed intervals down and merged
cover pieces as they unwound.  The iterative query walks \emph{up} from the two
boundary leaves, grabbing a node exactly when the boundary steps \emph{out} of
its parent on that side.  We phrase ranges half-open, \([l,r)\), which is what
makes the two-boundary bookkeeping clean.

\begin{lstlisting}[caption={half-open range fold on the 2n tree -- order-safe for any associative combine}, label={lst:c13-iter-range}]
long long query(int l, int r){      // fold over the half-open range [l, r)
    long long lres = ID, rres = ID; // TWO accumulators: keep left..right order
    for (l += n, r += n; l < r; l >>= 1, r >>= 1){
        if (l & 1) lres = mergeOp(lres, t[l++]); // left edge node: consumed L->R
        if (r & 1) rres = mergeOp(t[--r], rres); // right edge node: consumed R->L
    }
    return mergeOp(lres, rres);
}
// O(log n): at most two nodes folded per level.  For a COMMUTATIVE fold the
// two accumulators collapse into one variable and either merge order is fine.
\end{lstlisting}

\begin{proposition}{The \(2n\) tree supports point set and half-open range fold in \(\bigO{\log n}\)}{c13-iter}
For any associative \(\mathcal{M}\) with identity \(e\), \(\coderef{c13-iter}\)
builds a tree whose front-half node \(i\) holds \(\mathcal{M}\) applied to a fixed
dyadic block of leaves, \(\coderef{c13-iter}\)'s \code{update} restores that
invariant along the leaf's root path in \(\bigO{\log n}\), and
\(\coderef{c13-iter-range}\) returns \(\mathcal{M}\) over \([l,r)\) by folding
\(\le 2\lceil\log_2 n\rceil\) of those blocks.  If \(\mathcal{M}\) is also
\emph{commutative} (min, max, sum, gcd, xor, and, or), a single accumulator
suffices; otherwise the two-accumulator form above is required and is correct
because its left accumulator collects the cover blocks left-to-right and its
right accumulator collects them right-to-left, re-joined in order at the end.
\end{proposition}
\begin{proofbox}
The build loop fills indices \(n{-}1\) down to \(1\), so when \(t_i\) is written
its children \(t_{2i},t_{2i{+}1}\) (larger indices) already hold their blocks'
folds -- hence \(t_i\) becomes the fold of the union block; the invariant holds
by downward induction.  \code{update} rewrites one leaf and then re-runs
\(t_{p/2}=\mathcal{M}(t_p,t_{p\oplus 1})\) up the unique root path, exactly the
recursive \(\coderef{c13-update}\)'s re-folds, so the invariant is restored.  The
query maintains ``the answer is \(\mathcal{M}\) applied to (covered-left, the
still-open middle, covered-right): each halving step either leaves a boundary inside
its parent (nothing to grab, keep climbing) or steps to a boundary that is a
right child on the left edge (\code{l\&1}) / a left child on the right edge
(\code{r\&1}), and precisely then the sibling block lies wholly inside \([l,r)\)
and is folded in.  Blocks are folded in left-to-right order by the two
accumulators, and associativity lets them be re-associated into \(\mathcal{M}\)
of the whole range; commutativity makes the single-accumulator shortcut legal.
The number of grabbed blocks is at most two per level (one per boundary), so
\(\bigO{\log n}\).
\end{proofbox}

\begin{tracebox}[Half-open min query over positions 2..6 of the shared array]
Fold \(\min\) over \(A[1..5]\) (0-based), i.e.\ the half-open range
\([1,6)=\{2,4,5,1,1\}\), on the tree of \(\figref{c13-iter2n}\).  Start
\(\mathit{lres}{=}\mathit{rres}{=}e{=}{+}\infty\); \(\mathit{l}=1{+}8{=}9\),
\(\mathit{r}=6{+}8{=}14\).
\textbf{Step 1} \(\mathit{l}{=}9,\mathit{r}{=}14\): \(\mathit{l}\) odd \(\Rightarrow\)
\(\mathit{lres}{=}e\!\wedge\!t_9{=}\min(\infty,2){=}2\), \(\mathit{l}{=}10\);
\(\mathit{r}\) even \(\Rightarrow\) nothing; halve \(\mathit{l}{=}5,\mathit{r}{=}7\).
\textbf{Step 2} \(\mathit{l}{=}5,\mathit{r}{=}7\): \(\mathit{l}\) odd \(\Rightarrow\)
\(\mathit{lres}{=}\min(2,t_5{=}4){=}2\), \(\mathit{l}{=}6\); \(\mathit{r}\) odd
\(\Rightarrow\) \(\mathit{rres}{=}\min(t_6{=}1,\infty){=}1\), \(\mathit{r}{=}6\);
halve \(\mathit{l}{=}3,\mathit{r}{=}3\).  Loop ends (\(\mathit{l}<\mathit{r}\) fails).
Return \(\min(\mathit{lres},\mathit{rres}){=}\min(2,1){=}1\).  The cover grabbed
\(\{t_9,t_5,t_6\}=\{A_1\}\uplus\{A_2,A_3\}\uplus\{A_4,A_5\}=[1,6)\) -- three
disjoint blocks, minimum \(1\), cross-checking \(\secref{c13-core}\)'s recursive
answer on the same array.
\end{tracebox}

\subsection{Where Each Shape Wins}

The \(2n\) form is not strictly better -- it trades the recursive tree's
interval-awareness for array compactness.  Concretely: iterative \(2n\) wins on
memory (\(2n\) vs \(4n\)), constant factor (a flat loop, no call overhead, half
the cache lines of the padded tree), and code size; recursive wins wherever the
\emph{interval of a node is used} -- descending for a threshold
(\(\secref{c13-search}\)) needs \([l,r]\) at each node, and a non-commutative
range query is far less error-prone when the cover comes back in order from the
recursion instead of via the two-accumulator dance above.

\begin{comparisonbox}[Completing the boundary Ch.~18 drew (sum/min/max workload)]
Chapter~18's \(\secref{c14-range}\) settled \emph{expressiveness}: a BIT folds an
invertible or monotone accumulator, a segment tree folds \emph{any} associative
record.  With all five combines and both tree shapes now in hand, here are the
\emph{remaining} axes -- the ones Ch.~18 left as ``the code is there'':
\begin{center}\renewcommand{\arraystretch}{1.2}\footnotesize
\begin{tabular}{@{}p{0.115\linewidth} p{0.235\linewidth} p{0.255\linewidth} p{0.31\linewidth}@{}}
\toprule
 & \textbf{BIT} & \textbf{Rec. segtree} & \textbf{Iterative \(2n\) segtree}\\
\midrule
memory    & \(n{+}1\) & \(4n\) & \(2n\); pad to \(2^{\lceil\log_2 n\rceil{+}1}\) for ordered folds \\
point set & \(O(\log n)\) add only & \(O(\log n)\) set & \(O(\log n)\) set, one loop \\
range fold & \(O(\log n)\) \emph{prefix} only & \(O(\log n)\) any \([l,r]\) & \(O(\log n)\) any \([l,r)\) \\
combines  & invertible / monotone & \emph{any} assoc.\ record & any assoc.; commutative simplest \\
lines     & \(\sim\!8\) & \(\sim\!40\) & \(\sim\!15\) \\
best for  & sums/counts, tiny code & descend, non-commutative, lazy (Ch.20) & fixed length, hot inner loop \\
\bottomrule
\end{tabular}
\end{center}
\end{comparisonbox}

% ---------------------------------------------------------------
\subsection{Problem Lab: CF 514D R2D2 and Droid Army}
% ---------------------------------------------------------------

\solvedproblem{c13-r2d2}{R2D2 and Droid Army}{CF 514D}
  {https://codeforces.com/problemset/problem/514/D}{CF 2000}
  {one iterative 2n max-segtree per column + binary search on the window length; cost of a run = sum of per-column range maxima}

\begin{problemstatement}[CF 514D R2D2 and Droid Army]
An army of \(n\) droids is lined up in one row.  Each droid is described by
\(m\) integers \(a_1,\dots,a_m\), where \(a_i\) is the number of details of the
\(i\)-th type in its mechanism.  R2-D2 wants to destroy a run of
\emph{consecutive} droids of maximum length.  He has \(m\) weapons; the \(i\)-th
weapon destroys one detail of type \(i\) from every droid it is aimed at.  A
droid is destroyed once \emph{all} its details are gone, and R2-D2 may make at
most \(k\) shots in total.  Decide how many shots of each weapon type to make so
as to destroy the longest consecutive run.

\textbf{Input.} \(n,m,k\) with \(1\le n\le 10^5\), \(1\le m\le 5\),
\(0\le k\le 10^9\); then \(n\) lines of \(m\) values, each \(0\le a_i\le 10^8\).

\textbf{Output.} \(m\) space-separated integers: the shots of each type needed to
destroy a maximum-length consecutive run (any optimal run; fewer than \(k\) shots
is fine; the empty run gives all zeros).

\textbf{Examples.}
\begin{verbatim}
Input:            Output:      Input:      Output:
5 2 4             2 2          3 2 4       1 3
4 0                              1 2
1 2                              1 3
2 1                              2 2
0 2
1 3
\end{verbatim}
(First test destroys droids 2,3,4; second destroys droids 1,2.)
\end{problemstatement}

\begin{attemptgate}[Attempt it yourself first -- \lv{5}, give yourself one sitting]
To clear a \emph{run} of droids you must, for each type \(j\), fire at least
\(\max\) of \(a_j\) over the run -- the shots of type \(j\) is the run's
\emph{column maximum}.  So a run costs \(\sum_j(\text{column }j\text{ maximum})\),
and you want the longest run whose cost is \(\le k\).  Enumerating runs is
\(\bigO{n^2}\).
\peek{for a \emph{fixed} length \(L\), is it easy to test whether some window of
length \(L\) costs \(\le k\)?}
\nudge{``some window of length \(L\) works'' is monotone in \(L\): a sub-window's
column maxima are all \(\le\) the window's, so if length \(L\) is affordable so is
every \(L'<L\).  Binary-search the answer \(L\).}
\fullhint{For each of the \(m\) columns build an iterative \(2n\) max-segtree
(\(\coderef{c13-iter}\)).  \(\mathit{feasible}(L)\): slide a width-\(L\) window,
each column's contribution is that column's \(\text{query}(i,i{+}L)\)
(\(\coderef{c13-iter-range}\)); if the \(m\) maxima sum to \(\le k\) for some
start, \(L\) works.  Binary-search the largest feasible \(L\), then re-run
\(\mathit{feasible}\) once at that \(L\) to recover the printed vector.  Full
program below.}
\end{attemptgate}

\solstep{Understanding}
Fix a run of droids \([i, i{+}L)\).  To erase type \(j\) from every droid in it,
weapon \(j\) must fire \(W_j=\max_{p\in[i,i{+}L)} a_p^{(j)}\) shots -- that many
covers the droid with the most type-\(j\) details and harmlessly over-shoots the
rest, since a weapon hits the whole column.  The run costs \(\sum_j W_j\).  The
question ``is there a run of length \(L\) with cost \(\le k\)'' needs only the
\(m\) per-column range maxima of every width-\(L\) window, which is \(m\)
range-max structures -- the iterative \(2n\) tree, one per column (there are at
most \(m\le5\) of them, so building them all is cheap).

\solstep{Key Observations}
\begin{observation}
Feasibility is \textbf{monotone in \(L\)}: shrinking a window cannot raise any
column maximum, so \(\sum_j W_j\) never increases.  Therefore the set of feasible
lengths is a prefix \(0..L^\star\) and \(L^\star\) is found by binary search,
with \(\mathit{feasible}(\cdot)\) as the predicate.
\end{observation}
\begin{observation}
\(\max\) is associative, commutative, with identity \(-\infty\) (here values
\(\ge 0\), so \(0\) is a safe stand-in identity), so a single-accumulator
iterative query \(\coderef{c13-iter-range}\) is enough -- no order care needed.
\end{observation}
\begin{observation}
The answer vector is not a fresh computation: it is exactly the \(m\) column
maxima of \emph{any} optimal window.  Recover it by re-running the scan once at
the final \(L^\star\).
\end{observation}

\solstep{Think It Through With Me}
\begin{thinkbox}
\deadend{Two pointers with a multiset per column}
{The window ``cost \(\le k\)'' is not monotone under \emph{shrinking from the
left} in a usable way -- removing a droid may or may not drop the maximum, and you
would maintain \(m\) multisets to peek each column max.  Correct, but \(O(mn\log
n)\) with heavy constant and no cleaner than binary search; it also hides the
segment-tree lesson.}
\deadend{Re-scan every window from scratch inside the binary search}
{For a fixed \(L\), computing each window's column maxima naively is
\(O(L)\) per window \(\Rightarrow\) \(O(nL)\) per check \(\Rightarrow\)
\(O(n^2\log n)\) overall.  A range-max structure answers each window in
\(O(\log n)\), turning the check into \(O(mn\log n)\).}
\aha{One iterative max-segtree per column, binary-search the length, and read
each window's \(m\) maxima with \(m\) half-open queries; total
\(O(mn\log^2 n)\) with a tiny constant and about fifteen lines of tree code.}
\end{thinkbox}

\solstep{The Approach Ladder}
\begin{center}\renewcommand{\arraystretch}{1.25}
\begin{tabular}{@{}c p{0.22\linewidth} p{0.32\linewidth} l c@{}}
\toprule
\textbf{\#} & \textbf{Approach} & \textbf{Idea} & \textbf{Time} & \textbf{Verdict}\\
\midrule
1 & \lv{3}\ all windows & every \([i,j]\), column max on the fly & \(\bigO{m n^2}\) & \tle \\
2 & \lv{4}\ BS \(+\) naive scan & fixed \(L\), re-scan each window & \(\bigO{m n^2\log n}\) & \tle \\
3 & \lv{4}\ BS \(+\) multiset two-pointer & maintain \(m\) column maxes & \(\bigO{m n\log n}\) & passes, heavy \\
4 & \lv{5}\ \textbf{BS \(+\) iterative max-segtree} & \(m\) \(2n\) trees, half-open queries & \(\bigO{m n\log^2 n}\) & \ac \\
\bottomrule
\end{tabular}\end{center}

\solstep{Solution}
Store each column \(c\) as \(a_c[0..n{-}1]\).  Build \(\mathit{seg}[c]\) for all
\(c\).  \(\mathit{feasible}(L)\) returns false if \(L=0\) trivially (handled by the
search floor); for \(L\ge1\) loop start \(i=0..n{-}L\), sum the \(m\) values
\(\mathit{seg}[c].\mathit{query}(i,i{+}L)\), and succeed the moment that sum
\(\le k\), recording the vector.  Binary-search the largest feasible \(L\) over
\([0,n]\) (monotone by Observation~1), then print the vector from one more
\(\mathit{feasible}(L^\star)\) call -- or \(m\) zeros if \(L^\star=0\).

\solstep{Proof of Correctness}
\begin{proofbox}
\textbf{Cost lemma.} A run \([i,i{+}L)\) is destroyable within \(k\) shots iff
\(\sum_j \max_{p\in[i,i{+}L)}a_p^{(j)}\le k\): each type \(j\) needs exactly its
column maximum (enough for the worst droid, no less), and types are independent.
\(\mathit{feasible}(L)\) tests this for every \(L\)-window by reading the column
maxima from \(\mathit{seg}[c]\), which \propref{c13-iter} guarantees equals the
half-open range max.  \textbf{Monotonicity.} If \(L\) is feasible via window
\([i,i{+}L)\), then for any \(L'<L\) the sub-window \([i,i{+}L')\) has every
column maximum \(\le\) the length-\(L\) one, so its cost \(\le k\): feasible is
downward-closed, and binary search returns the largest feasible \(L^\star\).
\textbf{Answer.} \(\mathit{feasible}(L^\star)\) records the column maxima of a
concrete affordable length-\(L^\star\) window; printing them is a valid optimal
solution, and \(L^\star=0\) (no affordable single droid) prints the all-zero
vector, the empty run.
\end{proofbox}

\solstep{Implementation}
\begin{lstlisting}[caption={CF 514D full solution: per-column iterative max-segtrees + binary search on length}, label={lst:c13-514-soln}]
#include <bits/stdc++.h>
using namespace std;
const int MX = 100005;
int n, m; long long k;
int col[5][MX];                 // col[c][i] = type-c details of droid i, 0-based

struct MaxSeg {                 // iterative 2n max tree, values >= 0
    int N; vector<int> t;
    void build(int a[], int n){
        N = n; t.assign(2*N, 0);                // identity 0 (a_i >= 0)
        for (int i=0;i<n;i++) t[N+i] = a[i];
        for (int i=N-1;i>0;i--) t[i] = max(t[2*i], t[2*i+1]);
    }
    int query(int l, int r){                    // max over half-open [l, r)
        int res = 0;                            // commutative fold -> one acc
        for (l+=N, r+=N; l<r; l>>=1, r>>=1){
            if (l&1) res = max(res, t[l++]);
            if (r&1) res = max(res, t[--r]);
        }
        return res;
    }
} seg[5];

bool feasible(int L, int out[]){
    if (L == 0) return true;
    for (int i=0; i + L <= n; i++){
        long long cost = 0; int v[5];
        for (int c=0;c<m;c++){ v[c] = seg[c].query(i, i+L); cost += v[c]; }
        if (cost <= k){ for (int c=0;c<m;c++) out[c]=v[c]; return true; }
    }
    return false;
}
int main(){
    scanf("%d%d%lld", &n, &m, &k);
    for (int i=0;i<n;i++) for (int c=0;c<m;c++) scanf("%d", &col[c][i]);
    for (int c=0;c<m;c++) seg[c].build(col[c], n);
    int dummy[5];
    int lo = 0, hi = n, bestL = 0;              // largest feasible length
    while (lo <= hi){
        int mid = (lo + hi) / 2;
        if (feasible(mid, dummy)){ bestL = mid; lo = mid + 1; }
        else hi = mid - 1;
    }
    int out[5] = {0};
    feasible(bestL, out);                        // recover the answer vector
    for (int c=0;c<m;c++) printf("%d%c", out[c], c+1==m?'\n':' ');
}
\end{lstlisting}

\solstep{Trace the Official Samples}
\begin{tracebox}[CF 514D sample 1: n=5, m=2, k=4]
Binary search on \(L\in[0,5]\).  \(\mathit{feasible}(3)\): windows of length 3 --
\([0,3)\) gives column maxes \((\max 4,1,2,\ \max 0,2,1)=(4,2)\), cost \(6>4\);
\([1,4)\) gives \((\max 1,2,0,\ \max 2,1,2)=(2,2)\), cost \(4\le4\) \(\Rightarrow\)
success, record \((2,2)\).  \(\mathit{feasible}(4)\): \([0,4)=(4,2)\) cost 6,
\([1,5)=(2,3)\) cost 5, both \(>4\) \(\Rightarrow\) fail.  So \(L^\star=3\); one
more scan at \(L=3\) returns the affordable window \([1,4)\) (droids 2,3,4) with
vector \((2,2)\).  Output \texttt{2 2}.
\end{tracebox}
\begin{tracebox}[CF 514D sample 2: n=3, m=2, k=4]
\(\mathit{feasible}(2)\): \([0,2)=(\max 1,1,\max 2,3)=(1,3)\) cost \(4\le4\)
\(\Rightarrow\) success.  \(\mathit{feasible}(3)\): \([0,3)=(\max 1,1,2,\max
2,3,2)=(2,3)\) cost \(5>4\) \(\Rightarrow\) fail.  Hence \(L^\star=2\), the window
\([0,2)\) (droids 1,2) prints \((1,3)\).  Output \texttt{1 3}.
\end{tracebox}

\solstep{Complexity}
\begin{complexitybox}
Build: \(m\) trees of \(2n\) slots \(\Rightarrow\) \(\bigO{mn}\) time,
\(\bigO{mn}\) memory (\(m\le5\), \(n\le10^5\): about \(10^6\) \code{int}s,
\(\sim 4\) MB).  \(\mathit{feasible}(L)\) scans \(\le n\) windows, each folding
\(m\) half-open \(\bigO{\log n}\) queries \(\Rightarrow\) \(\bigO{mn\log n}\).
Binary search runs it \(\bigO{\log n}\) times: total \(\bigO{mn\log^2 n}\approx
5\cdot10^5\cdot 17^2 \approx 1.4\cdot10^8\) max/compare operations -- \textbf{AC}
at \(2\) s.  Swapping the per-window query for a width-\(L\) monotone deque would
drop one \(\log\); we keep the segment tree because it is the point of this
section and its \(2n\) constant is what makes the extra \(\log\) affordable.
\end{complexitybox}

\begin{hardeningbox}
\begin{itemize}[nosep,leftmargin=1.4em]
  \item \textbf{Print the maxima, not the length.}  The output is the \(m\) weapon
        counts of one optimal window; a common misread is to print \(L^\star\).
        A length-\(0\) answer is the all-zero vector, which the code emits for free.
  \item \textbf{Identity for a max tree.}  \(\coderef{c13-514-soln}\) uses \(0\)
        because \(a_i\ge0\); if negative values were allowed the identity would be
        \(-\infty\) and a careless \(0\) would silently over-report.
  \item \textbf{Half-open bounds.}  The query is \([i, i{+}L)\) for exactly \(L\)
        droids; writing \(i{+}L{+}1\) (a closed-range habit) reads one droid too
        many and breaks the monotone binary search's answer.
  \item \textbf{Recover the vector at \(L^\star\) only.}  Storing the vector during
        the search is a trap -- later feasible mids can be \emph{smaller} than
        \(L^\star\) and overwrite it with a shorter window's answer.
\end{itemize}
\end{hardeningbox}

% ---------------------------------------------------------------
\subsection{Common Pitfalls}
% ---------------------------------------------------------------

\begin{warningbox}[Trap 1: assuming the two accumulators can be merged]
For a \emph{commutative} fold the left and right accumulators collapse to one,
and the single-accumulator version of \(\coderef{c13-iter-range}\) is right.  For
an order-sensitive combine (max-subarray, brackets) the right-side blocks are
consumed right-to-left; folding them into one accumulator in the wrong order
\emph{reverses part of the range}.  Keep two accumulators, or use the recursive
\(\coderef{c13-query}\) which hands back the cover in order.
\end{warningbox}

\begin{warningbox}[Trap 2: non-power-of-two with a non-commutative fold]
The \(2n\) layout without padding is cleanest when \(n\) is a power of two; for
other \(n\) the implicit tree ``wraps'' a couple of ragged nodes, which is
harmless for commutative folds but scrambles intervals for ordered ones.  When
the fold is non-commutative and \(n\) is odd, pad to the next power of two with
the identity \(e\) (leaves beyond \(n\) equal \(e\)) and every node interval
becomes a proper dyadic block.
\end{warningbox}

\begin{warningbox}[Trap 3: forgetting the \(2n\) update climbs only siblings]
\(\coderef{c13-iter}\)'s \code{update} sets \(t_{p/2}=\mathcal{M}(t_p,
t_{p\oplus 1})\).  If you instead recompute a parent from two stale grandchildren
or forget the \code{p>1} guard, you either read index \(0\) or leave a node
un-propagated.  The recursion in \(\secref{c13-core}\) hides this because it
rebuilds from the returned child; the flat loop must name the sibling itself.
\end{warningbox}

\begin{warningbox}[Trap 4: binary search on a non-monotone predicate]
The \(2n\) tree made each \(\mathit{feasible}(L)\) cheap, but the \emph{outer}
binary search is valid only because \(\mathit{feasible}\) is monotone (shrinking a
window cannot grow a column maximum).  Before pairing a segtree with a
binary-search-on-answer, prove the monotonicity; here it is Observation~1, and
without it the search returns a wrong length even though every query was correct.
\end{warningbox}

% ---------------------------------------------------------------
\subsection{Practice Problems}
% ---------------------------------------------------------------

\begin{stuckbox}
\begin{itemize}[nosep,leftmargin=1.4em]
  \item Is your combine commutative?  If not, are you keeping two accumulators?
  \item Did you pad to a power of two (identity beyond \(n\)), or justify skipping
        the pad because the fold is commutative?
  \item In \code{update}, are you re-folding the parent from the child and its
        sibling \code{p\textasciicircum1}, not from stale values?
  \item For CF~514D: is the printed line the \(m\) maxima and not the length?
  \item Is your outer binary search's predicate actually monotone?
\end{itemize}
\end{stuckbox}

\begin{practiceset}
  \practice{Xenia and Bit Operations}{CF 339D}
    {https://codeforces.com/problemset/problem/339/D}
    {CF 1700 L3}
    {A perfect binary tree on \(n=2^m\) leaves where the merge \emph{alternates by
     level}: \(\mid\) (OR) one level above the leaves, then \(\oplus\) (XOR) above
     that, alternating to the root; point updates, print the root each time.  The
     fixed, complete shape is exactly the \(2n\) arena of \(\coderef{c13-iter},
     \coderef{c13-iter-range}\); the one wrinkle is that the fold is not a single
     \(\mathcal{M}\) but depends on the node's depth -- carry the parity (or build
     level by level in the loop) so each re-fold picks OR vs XOR.  Watch that the
     alternating merge is commutative, so a single accumulator still suffices.}
  \practice{Sereja and Brackets}{CF 380C}
    {https://codeforces.com/problemset/problem/380/C}
    {CF 2000 L4}
    {Given a bracket string, answer range queries: the length of the longest regular
     bracket \emph{subsequence} inside \([l,r]\).  A node is the triple (matched
     pairs, unmatched opens, unmatched closes) merged greedily; the merge is
     \textbf{non-commutative} (opens and closes do not swap roles), so this is the
     canonical stress test for \(\coderef{c13-iter-range}\)'s two accumulators -- a
     careless single-accumulator iterative query fails it while the recursive
     \(\coderef{c13-query}\) of \(\secref{c13-core}\) passes.  Reproduce the opener's
     max-subarray lesson (from \(\secref{c13-combine}\)) on a range, not the whole
     array.}
\end{practiceset}

\paragraph{Back-cites and forwards.}
\(\secref{c13-core}\) supplies the fold \(\mathcal{M}\), the identity \(e\), and
the \(4n\) form this section compacts to \(2n\); \(\secref{c13-combine}\) explains
the tuple records whose non-commutativity makes Trap~1 bite.  Chapter~18's
\(\secref{c14-range}\) drew the expressiveness line that this chapter's
comparison box completes.  Forward: \(\secref{c13-search}\)'s descend wants
node intervals and so leans recursive rather than \(2n\); \(\secref{c13-static}\)
pairs this fixed-length tree with offline range problems; and range \emph{updates}
-- the reason the \(2n\) loop still keeps a recursive cousin -- are Ch.~20's lazy
propagation (\(\chref{c15}\)).

% ---------------------------------------------------------------
\begin{recallbox}[Section 19.4 in 10 seconds]
A \textbf{complete} tree (fixed length) fits a flat \(\mathbf{2n}\) array: leaves
\(t[n..2n{-}1]\), internal \(t_i=\mathcal{M}(t_{2i},t_{2i+1})\) filled
\(i{=}n{-}1..\!1\); point set climbs by \(p{>}{>}1\) re-folding against sibling
\(p\oplus1\); range fold is the half-open \([l,r)\) two-accumulator walk
(\(\coderef{c13-iter-range}\), \(\propref{c13-iter}\)).  It wins on memory,
constant, and code; the recursive form wins for descends, non-commutative ranges,
and lazy.  CF~514D: one max-tree per column, binary-search the length on the
monotone ``cost \(\le k\)'' predicate, print the optimal window's \(m\) maxima.
\selftest{
  \item Why does the \(2n\) build need no interval arguments, and what does that
        cost you versus the recursive tree?
  \item When can the two accumulators of \(\coderef{c13-iter-range}\) collapse into
        one, and when must they stay separate?
  \item Give the identity for a max-tree here and the value-range fact that makes
        \(0\) safe.
  \item State the monotone predicate that makes CF~514D's binary search valid.
  \item Which is the honest answer to a length-\(0\) CF~514D case, and how does the
        code produce it?
}
\end{recallbox}
